\section{Background}
\label{sec:background}

\subsection{Two phases of LLM inference}
Autoregressive generation runs in two phases with very different hardware behavior.
In \emph{prefill}, all $L$ prompt tokens of each of the $B$ requests pass through the network at
once, so every linear layer is a GEMM with $M = B\cdot L$ rows. Its cost determines the
\emph{time to first token} (TTFT). In \emph{decode}, every step produces one new token per
request, so the same weights are multiplied by a matrix with only $M = B$ rows. Its cost per step
is the \emph{time per output token} (TPOT). The keys and values of all previous tokens are kept in
a KV cache whose size grows with $B\cdot(L+N)$ for $N$ generated tokens.

\subsection{A roofline view of quantization}
For a linear layer $y = xW^\top$ with $x\in\mathbb{R}^{M\times K}$ and $W\in\mathbb{R}^{N\times K}$
stored with $b$ bytes per weight, the arithmetic intensity is
\begin{equation}
  I(M) \;=\; \frac{2MNK}{b\,NK + 2MK + 2MN} \;\approx\; \frac{2M}{b}\quad (M \ll N, K).
  \label{eq:intensity}
\end{equation}
A kernel is memory-bound while $I(M)$ is below the GPU's ridge point
$P_{\text{peak}}/BW$ and compute-bound above it~\cite{williams2009roofline}. With BF16 weights
($b=2$), decode at small $B$ is deeply memory-bound: a step costs roughly the time needed to read
every weight once. Storing weights in 4 bits ($b\approx0.5$) cuts these bytes by about $4\times$,
which is the source of every decode-time speedup reported for weight-only quantization.
Three things work against it. (i) The weights must be dequantized back to 16 bits before the
tensor cores can use them, and this extra work is free only if a fused kernel hides it under the
memory stream. (ii) Weight-only formats (W4A16, W8A16) still compute in BF16, so once $I(M)$
reaches the ridge, at a $4\times$ smaller $M$ than BF16, the kernel is compute-bound and the
dequantization overhead is no longer hidden. (iii) Prefill is compute-bound from the start, so
weight-only formats cannot shorten TTFT. Formats that quantize the activations too (W8A8 in INT8
or FP8~\cite{micikevicius2022fp8}) also raise the compute ceiling, because the tensor cores run
8-bit math at twice the BF16 rate.

\subsection{Quantization methods studied}
\textbf{LLM.int8()}~\cite{dettmers2022llmint8} quantizes weights per output row and activations per
token to INT8. Activation feature columns that contain any value above a threshold (6.0) are
treated as outliers and multiplied in FP16 (mixed-precision decomposition).
\textbf{NF4}~\cite{dettmers2023qlora} stores weights in the 4-bit NormalFloat format with one
FP32 absmax scale per 64-weight block. Both are implemented in bitsandbytes and applied
to the BF16 checkpoint at load time through Hugging Face Transformers~\cite{hfbnb}.
\textbf{GPTQ}~\cite{frantar2023gptq} uses approximate second-order information to choose rounding
that minimizes each layer's output error. \textbf{AWQ}~\cite{lin2024awq} instead rescales the
weight channels that matter most for the activations before rounding. Both produce INT4 (or INT8)
weights with one FP16 scale per group of 128 weights, and are executed in vLLM~\cite{kwon2023vllm}
by the Marlin mixed-precision GEMM kernels~\cite{frantar2025marlin}. \textbf{FP8 W8A8} casts
the weights to E4M3 at load time with one scale per tensor and quantizes the activations
dynamically per token; vLLM runs it with CUTLASS FP8 GEMMs. SmoothQuant~\cite{xiao2023smoothquant} motivates the W8A8 design by moving
activation outliers into the weights; we use FP8 as its hardware-native counterpart on Blackwell.
